% Standalone source. Compile with pdflatex (twice) or tectonic.
% No external figures or bibliography files are required.
\documentclass[aps,pra,onecolumn,10pt,letterpaper,notitlepage,nofootinbib,amsmath,amssymb]{revtex4-2}
\usepackage[T1]{fontenc}
\usepackage{lmodern,mathtools,booktabs,microtype}
\usepackage{xcolor,hyperref}
\definecolor{navy}{HTML}{203B5D}
\hypersetup{colorlinks=true,linkcolor=navy,urlcolor=navy,citecolor=navy,
 pdftitle={Finite-depth QPM and Converter Efficiency},
 pdfsubject={Spatial overlap and a six-mode analytical transduction model}}
\newcommand{\kb}{k_{\mathrm B}}
\newcommand{\ER}{E_{\mathrm R}}
\newcommand{\kL}{k_{\mathrm L}}
\newcommand{\Om}{\Omega}
\newcommand{\etae}{\eta_{\mathrm{port}}}
\newcommand{\gadd}{\gamma_{m,\mathrm{add}}}

\begin{document}
\title{Finite-depth Quasi-Phase Matching (QPM) and Converter Efficiency}
\author{Abhishek V. Karve}
\affiliation{Department of Applied Physics, Stanford University}
\date{\today}
\maketitle
% {\LARGE\bfseries\color{black}Finite-depth QPM and converter efficiency\par}
% \vspace{4pt}
% {\color{muted}Spatial localization, collective-mode overlap, and an analytical six-mode model\par}
% \vspace{7pt}

An ensemble of $^{85}$Rb atoms couples to optical ($780\,\mathrm{nm}$) and microwave
cavities. A $480/481\,\mathrm{nm}$ blue control connects optical and Rydberg
excitations; $297\,\mathrm{nm}$ UV dressing enables the effective microwave coupling
from the prepared atomic state, completing the conversion loop~\cite{kumar}.
A separate $1188\,\mathrm{nm}$ lattice confines the atoms and provides the periodic
density for quasi-phase matching (QPM).

Finite site width reduces the coherent overlap $Q$ between the two atomic conversion
paths. This note derives the resulting efficiency, including cavity input and output
coupling, and gives a numerical example. The derivation is self-contained; the
companion Mathematica notebook implements and checks it, with a notation guide at
the end~\cite{model}.

In the level scheme of Ref.~\cite{kumar}, $g,e,r,f$ denote $5S_{1/2}$,
$5P_{3/2}$, $36S_{1/2}$, and $35P_{1/2}$, respectively. UV dressing mixes
$|g\rangle$ and $|f\rangle$ into the prepared reference state. Eliminating the unused
UV-dressed state leaves the effective microwave matrix element $g_m$ used below.
For optical-to-microwave conversion, the loop absorbs an optical and a blue photon
and emits a microwave and a UV photon. This sequence fixes the wavevector signs;
the reversed sequence describes microwave-to-optical conversion.

\section{What the finite-depth QPM factor means}
\label{sec:qpm}
For atoms $j=1,\ldots,N$, let $u_j,v_j$ be the complex excitation amplitudes in the
normalized optical/blue and UV/microwave collective profiles. They include relative
coupling magnitudes and phases; overall strengths enter separately in
section~\ref{sec:model}. Define
\begin{equation}
 Q=\left|\sum_j u_j^*v_j\right|,\qquad
 \sum_j|u_j|^2=\sum_j|v_j|^2=1,\qquad 0\leq Q\leq1.
 \label{eq:overlap}
\end{equation}
For matched, uniform envelopes, $Q=|\langle e^{i\Delta\mathbf k\cdot\mathbf r}\rangle|$,
averaged over the normalized atomic spatial distribution, with
$\Delta\mathbf k=\mathbf k_{780}+\mathbf k_{480}-\mathbf k_{\rm UV}-\mathbf k_\mu$.
Wavevectors include propagation directions; reversing the loop reverses
$\Delta\mathbf k$ without changing $Q$. Thus $Q$ is an amplitude overlap and $Q^2$
the squared useful-mode projection. Conversion probability also depends on coherent
exchange, atomic decay, and cavity ports.
Neither $Q$ nor $Q^2$ alone is a conversion probability.

\subsection{Localization within one site.}
An ideal retroreflected lattice has period $d=\lambda_{\mathrm L}/2$ and reciprocal
wavevector $K=2\kL$, where $\kL=2\pi/\lambda_{\mathrm L}$. For displacement $x$ from
a site minimum, barrier height $U$, and atomic mass $m$, the harmonic approximation
gives the recoil energy $\ER$ and trap angular frequency $\omega_t$~\cite{bloch}:
\begin{equation}
 V(x)=U\sin^2(\kL x)\simeq U\kL^2x^2,\quad
 \ER=\frac{\hbar^2\kL^2}{2m},\quad
 \omega_t=\frac{2\sqrt{U\ER}}{\hbar}.
 \label{eq:trap}
\end{equation}
With reduced Planck constant $\hbar$, Boltzmann constant $\kb$, and temperature $T$,
the thermal position variance and density Fourier harmonics are
\begin{align}
 \sigma_x^2&=\frac{\hbar}{2m\omega_t}
       \coth\!\left(\frac{\hbar\omega_t}{2\kb T}\right),\nonumber\\[-1pt]
 b_n&=\left|\left\langle e^{inKx}\right\rangle\right|
     =e^{-(nK)^2\sigma_x^2/2}
     =\exp\!\left[-n^2\sqrt{\frac{\ER}{U}}
           \coth\!\left(\frac{\sqrt{U\ER}}{\kb T}\right)\right].
 \label{eq:bn}
\end{align}
The integer $n$ is a spatial harmonic order, not a vibrational band. At $T=0$,
$\coth\to1$, but zero-point motion keeps $b_n<1$ for finite $U$ and $n\ne0$.
Deeper trapping drives $b_n\to1$ at fixed $T$. This approximation requires
localization near the minimum; occupation near the barrier requires band wavefunctions.

\subsection{From one axis to the total overlap.}
Decompose position into the site centre $\mathbf R$ and the in-site displacement.
Exact matching, $\Delta\mathbf k=\sum_i n_i\mathbf K_i$ with integer $n_i$, makes
equivalent site centres contribute in phase. Finite in-site widths still reduce
the overlap; for independent confined directions,
\begin{equation}
 Q=\prod_i b_{n_i,i};\qquad
 Q=b_{2,x}b_{2,y}\quad\hbox{for two second-order QPM axes}.
 \label{eq:qproduct}
\end{equation}
The last expression is the working geometry used below: co-propagating 780 and
480/481 nm fields, perpendicular UV, and suitable lattice axes. For an illustrative
orthogonal-axis layout, take
$\mathbf K_x=K\hat{\mathbf x}$, $\mathbf K_y=K\hat{\mathbf y}$,
$\mathbf k_{780},\mathbf k_{\rm blue}\parallel+\hat{\mathbf x}$, and
$\mathbf k_{\rm UV}\parallel+\hat{\mathbf y}$. The target harmonic is then
$(n_x,n_y)=(2,-2)$; $b_{-2}=b_2$. The component matching conditions are
\begin{equation}
 \begin{aligned}
 k_{780}+k_{\rm blue}-k_{\mu,x}&=2K,\\
 -k_{\rm UV}-k_{\mu,y}&=-2K,\qquad k_{\mu,z}=0.
 \end{aligned}
 \label{eq:geometry-components}
\end{equation}
This is a \emph{conditional geometry}, not fixed by the nominal wavelengths.
For example, neglecting $\mathbf k_\mu$ and taking the labels literally gives
$\delta k_x=-0.01017\,\mathrm{rad}/\mu\mathrm m$ for 480 nm and
$-0.03738\,\mathrm{rad}/\mu\mathrm m$ for 481 nm, with 780 nm and 1188 nm
held fixed. These correspond to about 1.02 and 3.74 radians across $100\,\mu\mathrm m$.
Actual wavelengths, beam angles, and reciprocal vectors must satisfy the vector
condition; the numerical overlap below assumes that they do.

A residual mismatch $\delta\mathbf k$ produces phase cancellation across occupied
site centres through $|\langle e^{i\delta\mathbf k\cdot\mathbf R}\rangle|$.
For Gaussian site-centre widths $L_i$, this factor is
$\exp[-\sum_i(\delta k_iL_i)^2/2]$, so $|\delta k_i|L_i\ll1$ is a useful
tolerance criterion. The in-site Fourier factor must also be evaluated at the
actual mismatch. Exact site-centre matching aligns the phases from different wells;
finite width still averages different phases within each well. Mode envelopes
must be checked separately.

\section{Analytical converter model with the orthogonal modes}
\label{sec:model}
\subsection{Problem definition and collective-mode construction.}
The system contains an optical cavity mode $\hat a$, a microwave cavity mode
$\hat b$, and an ensemble with optically excited and Rydberg excitations. The optical
cavity couples to one collective atomic profile with strength $g_o$; the UV-dressed
microwave interaction couples to another with strength $g_m$. The blue field connects
the optical and Rydberg excitations with matrix element $\Om$. The profiles need not
coincide: their normalized overlap is $Q$ from Eq.~\eqref{eq:overlap}.

We drive only the optical input and seek the coherent microwave output in the
linear, weak-excitation regime. Efficiency is the ratio of output microwave photon
flux to input optical photon flux; control fields carry the photon-energy difference.
This is a photon-number efficiency, not an energy-efficiency ratio.
All analytical couplings, detunings, and linewidths are angular frequencies;
section~\ref{sec:example} gives numerical values divided by $2\pi$.

\subsubsection{Atomic variables and assumptions.}
Let $\hat e_j^\dagger$ and $\hat r_j^\dagger$ create optical and Rydberg excitations
on atom $j$ relative to the prepared UV-dressed reference state. The optical
excitation is its coherence with $|e\rangle$; the Rydberg excitation is its
coherence with $|r\rangle$, addressed by the blue and UV-dressed microwave paths. When signal
excitations are few compared with the prepared atoms, depletion and saturation
are negligible and these operators are approximately bosonic:
\begin{equation}
 [\hat e_j,\hat e_\ell^\dagger]\simeq
 [\hat r_j,\hat r_\ell^\dagger]\simeq\delta_{j\ell},\qquad
 [\hat e_j,\hat r_\ell^\dagger]\simeq0.
 \label{eq:atomic-bosons}
\end{equation}
Assume frozen positions, uniform blue-control magnitude, identical atomic detunings
and independent decay rates. UV dressing is included in $g_m$ and its profile.
Absorbing the blue phase into $\hat r_j$ makes $\Om$ real and uniform. It is the
Hamiltonian matrix element below; convert any Rabi-frequency convention containing
an extra factor of $1/2$ before substitution:
\begin{equation}
 \frac{\hat H_B}{\hbar}
 =-\Om\sum_j\bigl(\hat e_j^\dagger\hat r_j+
                         \hat r_j^\dagger\hat e_j\bigr).
 \label{eq:uniform-blue}
\end{equation}

\subsubsection{The two coupling profiles.}
In this phase convention, the optical cavity selects $\mathbf u$ and the UV-dressed
microwave interaction selects $\mathbf v$, as in section~\ref{sec:qpm}:
\begin{equation}
 \mathbf u^\dagger\mathbf u=\mathbf v^\dagger\mathbf v=1,\qquad
 \mathcal Q=\mathbf u^\dagger\mathbf v=\sum_j u_j^*v_j,\qquad
 Q=|\mathcal Q|.
 \label{eq:complex-overlap}
\end{equation}
Absorb the global phase of $\mathbf v$ into $\hat b$ to choose
$\mathbf u^\dagger\mathbf v=Q\geq0$. With $s=\sqrt{1-Q^2}$,
\begin{equation}
 \boxed{\mathbf v=Q\mathbf u+s\mathbf w,}\qquad
 \mathbf w=\frac{\mathbf v-Q\mathbf u}{s},\qquad
 \mathbf u^\dagger\mathbf w=0,\quad\mathbf w^\dagger\mathbf w=1.
 \label{eq:profile-decomposition}
\end{equation}
The formula for $\mathbf w$ applies for $Q<1$. At $Q=1$, any normalized orthogonal
profile suffices because its microwave coupling vanishes ($s=0$).

\subsubsection{Collective operators and closure.}
Define four collective operators (uppercase letters distinguish operators from
their mean amplitudes):
\begin{equation}
 \begin{aligned}
 \hat P^\dagger&=\sum_j u_j\hat e_j^\dagger,&
 \hat R^\dagger&=\sum_j u_j\hat r_j^\dagger,\\
 \hat P_\perp^\dagger&=\sum_j w_j\hat e_j^\dagger,&
 \hat R_\perp^\dagger&=\sum_j w_j\hat r_j^\dagger.
 \end{aligned}
 \label{eq:collective-operators}
\end{equation}
In the weak-excitation limit, orthonormality gives unit commutators within each mode;
distinct modes commute. The coupling chain is
\begin{equation}
 a
 \xleftrightarrow{g_o} P
 \xleftrightarrow{\Om} R
 \xleftrightarrow{Qg_m} b
 \xleftrightarrow{s g_m} R_\perp
 \xleftrightarrow{\Om} P_\perp .
 \label{eq:coupling-chain}
\end{equation}
Here the labels denote modes, with hats suppressed. The orthogonal pair describes
a different spatial pattern of excitation of the same atoms. It has no direct
optical-cavity coupling but remains coupled to the microwave cavity and blue field.
Orthogonality alone is not irreversible loss; atomic decay makes this channel
dissipative. The two cavity coupling profiles span at most a
two-dimensional atomic subspace, preserved by uniform control, detunings, and damping.
Consequently, one additional optical/Rydberg pair suffices: the response closes on
four atomic and two cavity modes. All other orthogonal atomic modes are initially
unexcited and have zero mean drive.

\subsubsection{Hamiltonian and conventions.}
The collective matrix elements $g_o,g_m$ include each profile's total coupling
magnitude and are held fixed as $Q$ varies. They are collective couplings, not
single-atom matrix elements. The baseline parameter convention is~\cite{model}
\begin{equation}
 g_o=\frac{G}{\sqrt2},\qquad g_m=\frac{G_\mu}{\sqrt2}.
 \label{eq:coupling-convention}
\end{equation}
Thus $G,G_\mu$ rescale the same collective couplings. These factors reproduce the
notebook convention; they are not an additional finite-depth correction.
For experimental calibration, $g_o,g_m$ should be identified with the measured
collective matrix elements in Eq.~\eqref{eq:interaction-hamiltonian}.
With zero optical-cavity
detuning, the rotating-frame Hamiltonian fixes the signs of the optical atomic,
Rydberg, and microwave detunings $\Delta_e,\Delta_r,\Delta_\mu$:
\begin{align}
 \frac{\hat H_0}{\hbar}
 &=-\Delta_\mu\hat b^\dagger\hat b
   -\Delta_e(\hat P^\dagger\hat P+\hat P_\perp^\dagger\hat P_\perp)
   -\Delta_r(\hat R^\dagger\hat R+\hat R_\perp^\dagger\hat R_\perp),
   \label{eq:free-hamiltonian}\\
 \frac{\hat H_{\rm int}}{\hbar}
 &=-\Bigl[g_o\hat a^\dagger\hat P
   +\Om(\hat P^\dagger\hat R+\hat P_\perp^\dagger\hat R_\perp)
   +g_m\hat b^\dagger(Q\hat R+s\hat R_\perp)+\mathrm{h.c.}\Bigr].
   \label{eq:interaction-hamiltonian}
\end{align}
Use $\dot{\hat c}=(i/\hbar)[\hat H,\hat c]$ and Fourier amplitudes
$e^{+i\omega t}$, with spectral offset $\omega$. For optical input and useful
microwave output coupling rates $\kappa_{o,\mathrm{ex}},\kappa_{m,\mathrm{ex}}$,
the optical drive is
\begin{equation}
 \frac{\hat H_{\rm d}}{\hbar}
 =-\sqrt{\kappa_{o,\mathrm{ex}}}\,
   \bigl[\hat a^\dagger a_{\rm in}(t)+a_{\rm in}^*(t)\hat a\bigr].
 \label{eq:drive-hamiltonian}
\end{equation}
Here $|a_{\rm in}|^2$ is incident photon flux. Write
$a=\langle\hat a\rangle$, $b=\langle\hat b\rangle$, $p=\langle\hat P\rangle$,
$r=\langle\hat R\rangle$, $p_\perp=\langle\hat P_\perp\rangle$, and
$r_\perp=\langle\hat R_\perp\rangle$.
With no microwave input, $b_{\rm out}=\sqrt{\kappa_{m,\mathrm{ex}}}\,b$ up to a
port phase, and $S_{mo}=b_{\rm out}/a_{\rm in}$. Reservoir noise has zero mean;
this response calculation does not characterize added noise.

\subsection{Response functions and linewidth conventions.}
Let $\kappa_o,\kappa_m$ be the bare cavity energy decay rates, including all ports
and internal cavity loss. Thus $\kappa_{o,\mathrm{ex}}\leq\kappa_o$ and
$\kappa_{m,\mathrm{ex}}\leq\kappa_m$. The atomic linewidth parameters
$\Gamma,\gamma_r$ give optical and Rydberg amplitude decay rates
$\Gamma/2,\gamma_r/2$, respectively.

Introduce an independent added microwave energy-loss rate $\gadd\geq0$ for any
photon-loss reservoir not already included in $\kappa_m$ or the explicit atomic
dynamics. The baseline notebooks set $\gadd=\gamma_r$~\cite{model}, but do not
identify a microscopic reservoir for that choice. We retain this equality only
for the legacy numerical benchmark. The linearized cavity equation in
Ref.~\cite{kumarSI}, Sec.~B.3, instead has bare damping $\kappa_m/2$;
it corresponds to $\gadd=0$ here. Atomic loss mediated by the retained Rydberg
and optical modes is already accounted for by their response and must not be
counted again as $\gadd$. A physical use of nonzero $\gadd$ therefore requires
independent justification. Including damping, the inverse response functions are
\begin{equation}
 \begin{aligned}
 D_a&=\kappa_o/2+i\omega,&
 D_e&=\Gamma/2+i(\omega-\Delta_e),\\
 D_r&=\gamma_r/2+i(\omega-\Delta_r),&
 D_b&=(\kappa_m+\gadd)/2+i(\omega-\Delta_\mu).
 \end{aligned}
 \label{eq:D}
\end{equation}
Real parts give amplitude damping; imaginary parts give detuning. Orthogonal-mode
feedback is derived separately below.

\subsection{Six coupled equations.}
With these conventions,
\begin{equation}
 \begin{aligned}
 D_a a-i g_o p &=i\sqrt{\kappa_{o,\mathrm{ex}}}\,a_{\rm in},\\
 D_e p-i g_o a-i\Om r&=0,\\
 D_r r-i\Om p-i g_m Q b&=0,\\
 D_b b-i g_m\bigl(Qr+\sqrt{1-Q^2}\,r_\perp\bigr)&=0,\\
 D_r r_\perp-i\Om p_\perp-i g_m\sqrt{1-Q^2}\,b&=0,\\
 D_e p_\perp-i\Om r_\perp&=0.
 \end{aligned}
 \label{eq:six}
\end{equation}
The microwave coupling splits into useful amplitude $Qg_m$ and orthogonal amplitude
$g_m\sqrt{1-Q^2}$, retaining total strength $g_m$.

\subsection{Eliminating the orthogonal pair.}
Substitute $p_\perp=i\Om r_\perp/D_e$ from the last equation into the preceding
one. With $F=D_eD_r+\Om^2$,
\begin{equation}
 r_\perp=\frac{i g_m\sqrt{1-Q^2}\,D_e}{F}\,b,
 \qquad
 \boxed{\Sigma_\perp(\omega)=\frac{g_m^2(1-Q^2)D_e}{F}}.
 \label{eq:sigma}
\end{equation}
The microwave inverse response becomes $D_b+\Sigma_\perp$. The eliminated pair
adds amplitude damping $\operatorname{Re}\Sigma_\perp$ (full linewidth
$2\operatorname{Re}\Sigma_\perp$) and dispersion $\operatorname{Im}\Sigma_\perp$.
Despite lacking direct optical-cavity coupling, it is microwave-driven, feeds back
on that field, and dissipates through atomic decay. Replacing $G_\mu$ by $QG_\mu$
alone discards this feedback. Neither this damping nor the baseline extra damping
belongs in the useful output-port rate.

\subsection{Closed transfer function.}
With $B_\perp=D_b+\Sigma_\perp$, eliminating the remaining atomic amplitudes gives
\begin{equation}
 S_{mo}=\frac{\sqrt{\kappa_{o,\mathrm{ex}}\kappa_{m,\mathrm{ex}}}\,
                    g_og_m\Om Q}
 {B_\perp(D_aF+g_o^2D_r)+g_m^2Q^2(D_aD_e+g_o^2)}.
 \label{eq:transfer-eliminated}
\end{equation}
Set $Y=D_aF+g_o^2D_r$, $Z=D_bF+g_m^2D_e$. Substituting $\Sigma_\perp$ and
$F=D_eD_r+\Om^2$ gives, up to the port phase,
\begin{equation}
 \boxed{
 S_{mo}(\omega)=
 \frac{\sqrt{\kappa_{o,\mathrm{ex}}\kappa_{m,\mathrm{ex}}}\,
       g_og_m\Om QF}
      {YZ+Q^2g_o^2g_m^2\Om^2},\qquad
 \eta(\omega)=|S_{mo}(\omega)|^2 .}
 \label{eq:transfer}
\end{equation}
Because $Q$ enters both numerator and denominator, efficiency cannot generally be
scaled by $Q^2$. At $Q=1$ the orthogonal pair decouples, recovering the four-mode
baseline; at $Q=0$ coherent conversion vanishes.

\subsection{Effective cavity coupling and matching.}
Eliminating all four atomic amplitudes at once exposes the physical structure:
\begin{equation}
 \begin{pmatrix} A&iJ\\ iJ&B\end{pmatrix}
 \begin{pmatrix}a\\b\end{pmatrix}
 =\begin{pmatrix}i\sqrt{\kappa_{o,\mathrm{ex}}}\,a_{\rm in}\\0\end{pmatrix},
 \qquad
 A=D_a+\frac{g_o^2D_r}{F},\quad
 B=D_b+\frac{g_m^2D_e}{F},\quad
 J=\frac{Qg_og_m\Om}{F}.
 \label{eq:effective-cavities}
\end{equation}
Here $A$ and $B$ are the fully dressed inverse cavity responses; $B$ differs from
the intermediate $B_\perp$. The response is
$S_{mo}=\sqrt{\kappa_{o,\mathrm{ex}}\kappa_{m,\mathrm{ex}}}\,J/(AB+J^2)$,
equivalent to Eq.~\eqref{eq:transfer}. This frequency-domain elimination is exact
within Eq.~\eqref{eq:six}; no assumption of fast atomic relaxation is needed.
Away from resonance, $A,B,J$ are complex, frequency-dependent susceptibilities.

At fixed $g_o,g_m,\Om$ and linewidths, the diagonal atomic loading of each cavity
is independent of $Q$, while the cross-cavity coupling $J$ is proportional to $Q$.
The microwave mode still couples to its entire normalized atomic profile:
its useful and orthogonal contributions to the diagonal loading add as
$Q^2+(1-Q^2)=1$. Reducing $g_m$ to $Qg_m$ would reduce both the cross coupling
and that total loading, explaining why it gives a different response.

On resonance, $A,B>0$ and $J$ are real. Define
\begin{equation}
 \mathcal C=\frac{J^2}{AB},\qquad
 \boxed{\eta(0)=
 \frac{\kappa_{o,\mathrm{ex}}}{2A}
 \frac{\kappa_{m,\mathrm{ex}}}{2B}
 \frac{4\mathcal C}{(1+\mathcal C)^2}.}
 \label{eq:matching}
\end{equation}
The first two factors compare useful port escape with the dressed damping.
The final factor describes coupling relative to damping and reaches unity at
$\mathcal C=1$: the two damped responses are matched. This condition maximizes
conversion at fixed $A,B$, but need not eliminate reflection when internal losses
are present. For $\mathcal C<1$, weaker overlap moves the system farther from
matching and lowers conversion. For $\mathcal C>1$, reducing $Q$ can instead
improve conversion at the fixed resonant frequency by bringing $\mathcal C$
toward one. Thus $Q$ is a mode overlap, not a universal multiplicative efficiency
penalty; a loss budget also requires the reflected field.

\section{On-resonance efficiency and the numerical value of \texorpdfstring{$Q$}{Q}}
\label{sec:example}
At $\omega=\Delta_e=\Delta_r=\Delta_\mu=0$, define cooperativities comparing
coherent coupling to damping:
\begin{equation}
 C_o=\frac{2G^2}{\kappa_o\Gamma},\qquad
 C_m=\frac{2G_\mu^2}{(\kappa_m+\gadd)\gamma_r},\qquad
 C_B=\frac{4\Om^2}{\Gamma\gamma_r}.
 \label{eq:coops}
\end{equation}
Let $\eta_o=\kappa_{o,\mathrm{ex}}/\kappa_o$ and
$\eta_m=\kappa_{m,\mathrm{ex}}/\kappa_m$. Equation~\eqref{eq:transfer} becomes
\begin{equation}
 \boxed{
 \eta(0)=\eta_o\eta_m\,\frac{\kappa_m}{\kappa_m+\gadd}\,
 \frac{4Q^2C_oC_mC_B(1+C_B)^2}
 {\left[(1+C_B+C_o)(1+C_B+C_m)+Q^2C_oC_mC_B\right]^2}.}
 \label{eq:eta0}
\end{equation}
In these variables, the matching parameter is
\begin{equation}
 \mathcal C=\frac{Q^2C_oC_mC_B}
 {(1+C_B+C_o)(1+C_B+C_m)}.
 \label{eq:matching-coops}
\end{equation}
Port fractions multiply the efficiency at fixed total linewidths. Changes to
mirror or microwave-port coupling that alter linewidths require updated
cooperativities. Propagation and detection transmissions further multiply this
port-to-port result to give end-to-end efficiency.

\subsection{Lattice example.}
Take $^{85}$Rb, $\lambda_{\mathrm L}=1188\,\mathrm{nm}$, waist
$w=200\,\mu\mathrm m$, and 100\% retroreflection. Adopt the specified calibration
$3\,\mathrm W$ per beam per axis $\leftrightarrow U/\kb=15\,\mu\mathrm K$;
no polarizability calculation is implied. With $\ER/\kb=0.0799004\,\mu\mathrm K$,
$T=0$, and two equal, exactly matched second-order axes,
\begin{equation}
 b_2=e^{-4\sqrt{\ER/U}}=0.746816,\qquad
 \boxed{Q=b_2^2=0.557734},\qquad Q^2=0.311067.
 \label{eq:qnumber}
\end{equation}
This is a harmonic ground-state benchmark. The relevant length and energy scales are
\begin{equation}
 \begin{aligned}
 U/\ER&=187.7,& \sigma_x&=36.12\,\mathrm{nm},& 2K\sigma_x&=0.7641,\\
 \omega_t/(2\pi)&=45.62\,\mathrm{kHz},&
 \hbar\omega_t/\kb&=2.190\,\mu\mathrm K,& 2\pi/\omega_t&=21.92\,\mu\mathrm s.
 \end{aligned}
 \label{eq:lattice-scales}
\end{equation}
Even though the lattice is deep in recoil units, each second-order axis has
an rms mismatch phase of $0.7641$ radians across a site. Averaging these phases
gives $b_2=e^{-(0.7641)^2/2}$; multiplying two independent axes gives $Q=0.557734$.
The phase spread $nK\sigma_x$, rather than depth alone, controls the overlap and
explains the strong $n^2$ sensitivity. Ground-state occupation requires temperature
small compared with $\hbar\omega_t/\kb$. For example, Eq.~\eqref{eq:bn} gives
$Q=0.48137$ at $T=1\,\mu\mathrm K$, with the same depth and geometry.

The equilibrium density supplies an equal-time spatial average, including
zero-point width. Replacing the atomic motion by this static profile additionally
requires little motional evolution over the converter response time
$\tau_{\rm resp}$: $\omega_t\tau_{\rm resp}\ll1$, with motional sidebands
unresolved by the relevant response. This is distinct from requiring small
$nK\sigma_x$; the Gaussian phase average above does not expand in that quantity.
The timescale check for the converter benchmark is given below. The model predicts
the coherent response and does not calculate the inelastic motional spectrum.

\subsection{Converter example.}
Use the collective couplings and linewidths
\begin{equation}
 \begin{aligned}
 \frac{(G,G_\mu,\Om)}{2\pi}
   &=(15.8474565,\,2.09170145,\,5.52497339)\,\mathrm{MHz},\\
 \frac{(\kappa_o,\kappa_m,\Gamma,\gamma_r)}{2\pi}
   &=(10,\,5,\,6,\,0.1)\,\mathrm{MHz}.
 \end{aligned}
 \label{eq:benchmark}
\end{equation}
All detunings are zero, and $\gadd=\gamma_r$ reproduces the notebook convention.
Consistently dividing every rate and frequency by $2\pi$
leaves the efficiency unchanged, allowing direct use of these MHz values.
For unit port fractions and the finite-depth $Q$ of Eq.~\eqref{eq:qnumber}:
\begin{center}
\begin{tabular}{@{}lc@{}}
\toprule
Model & On-resonance efficiency\\
\midrule
Matched profiles, $Q=1$ & $82.0957\%$\\
Naive probability scaling, $Q^2\eta(Q=1)$ & $25.5373\%$\\
Only $G_\mu\mapsto QG_\mu$; orthogonal pair omitted & $51.5929\%$\\
Six modes with $Q=0.557734$ & $\mathbf{47.0824\%}$\\
\bottomrule
\end{tabular}
\end{center}
Here $\Sigma_\perp(0)/(2\pi)=0.147394\,\mathrm{MHz}$ and the additional full
linewidth is $2\operatorname{Re}\Sigma_\perp(0)/(2\pi)=0.294787\,\mathrm{MHz}$.
Port fractions $\eta_o=\eta_m=0.98$ give $45.21798\%$ at unchanged total linewidths,
before further path losses.

For this benchmark, $A/(2\pi)=5.20468\,\mathrm{MHz}$ and
$B/(2\pi)=2.76394\,\mathrm{MHz}$ on resonance. The matching parameter falls from
$\mathcal C=0.61947$ at $Q=1$ to $0.19270$ at the finite-depth overlap.
Both cases are below one, so reduced overlap worsens matching. At $T=1\,\mu\mathrm K$,
the smaller overlap quoted above would give $\eta(0)=38.15\%$ with all converter
parameters unchanged. Setting only $\gadd=0$ instead gives $48.54\%$ at $T=0$;
this isolates the effect of the unassigned extra microwave-loss assumption.

The six-mode homogeneous drift matrix provides a response-time check. At the
finite-depth benchmark its slowest amplitude decay rate is
$2\pi\times1.4533\,\mathrm{MHz}$, corresponding to
$\tau_{\rm resp}=0.1095\,\mu\mathrm s$. Thus
$\omega_t\tau_{\rm resp}=0.0314\ll1$, and the $45.62\,\mathrm{kHz}$ trap
spacing is small compared with the response relaxation scale. This supports
the frozen-motion approximation for this benchmark; parameters producing a much
narrower resonance require a new check and possibly a dynamical motional model.

\subsection{Where the unconverted input goes.}
For the drive phase convention in Eq.~\eqref{eq:six}, the optical reflected field is
$a_{\rm out}=a_{\rm in}+i\sqrt{\kappa_{o,\mathrm{ex}}}\,a$.
With incident flux $\Phi=|a_{\rm in}|^2$, define
\begin{equation}
 \begin{aligned}
 \mathcal R_o&=\left|1+i\sqrt{\kappa_{o,\mathrm{ex}}}\,a/a_{\rm in}\right|^2,\\
 \mathcal L_{\rm at}&=\frac{\Gamma(|p|^2+|p_\perp|^2)
                       +\gamma_r(|r|^2+|r_\perp|^2)}{\Phi},\\
 \mathcal L_o&=\frac{(\kappa_o-\kappa_{o,\mathrm{ex}})|a|^2}{\Phi},\qquad
 \mathcal L_m=\frac{(\kappa_m-\kappa_{m,\mathrm{ex}}+\gadd)|b|^2}{\Phi}.
 \end{aligned}
 \label{eq:flux-budget}
\end{equation}
The passive linear model satisfies
$\mathcal R_o+\eta+\mathcal L_{\rm at}+\mathcal L_o+\mathcal L_m=1$.
For unit port fractions, $\mathcal L_o=0$ and only $\gadd$ contributes to
$\mathcal L_m$. The benchmark budget, in percent of incident photon flux, is
\begin{center}
\begin{tabular}{@{}l@{\hspace{2em}}r@{\hspace{1.5em}}r@{}}
\toprule
Channel & $Q=1$ & $Q=0.557734$\\
\midrule
Microwave output & 82.0957 & 47.0824\\
Optical reflection & 3.4748 & 37.3236\\
Total atomic dissipation & 12.7875 & 14.6524\\
Added microwave dissipation & 1.6419 & 0.9416\\
\midrule
Total (before rounding) & 100 & 100\\
\bottomrule
\end{tabular}
\end{center}
The orthogonal pair accounts for $2.7759\%$ of incident flux, included in total
atomic dissipation. The conversion falls by $35.0133$ percentage points, while
reflection rises by $33.8487$ points and total atomic dissipation by only $1.8648$
points. Most of the conversion reduction therefore appears as increased reflection.
The orthogonal channel changes the coupled response as well as dissipating:
its direct loss is not equal to the difference between the four- and six-mode
efficiencies. This is why neither projection loss nor added linewidth alone
explains the output.

\subsection{Scope and checks.}
The closure requires weak excitation, frozen positions, uniform control and atomic
linewidths, and one effective optical profile. Anharmonic motion, resolved sidebands,
differential light shifts, or further independent standing-wave channels require
an extended susceptibility. Mathematica checks reproduce the six-equation solution
and $Q=1$ limit. Direct matrix solutions also verify the effective-cavity reduction,
the numerical tables, and photon-flux conservation. These are checks of algebra
and numerical implementation; they do not establish exact lattice matching or a
microscopic origin for $\gadd$. Numerical predictions remain conditional on the
stated geometry and model assumptions.

\subsection{Using the companion Mathematica notebook.}
\texttt{Abhishek\_Converter\_QPM.nb} solves the six equations, eliminates the
orthogonal pair, checks $Q=1$, and evaluates the example. Its
$\kappa,\kappa_\mu$ are $\kappa_o,\kappa_m$ here; other physical symbols are
unchanged. The calculator's $G_{\rm opt},G_{\rm mm}$ are $G,G_\mu$.
Set the auxiliary parameter $\epsilon=1$. Both notebooks specialize to
$\gadd=\gamma_r$; the independent $\gadd$ in this note generalizes their
microwave damping. All quoted notebook comparisons use that specialization.

The notebook's forcing parameter \texttt{Ein} remains in its amplitude
$\texttt{ConverterEfficQ}=\sqrt{\kappa_o\kappa_m}\,b$.
The input-normalized response, up to the port phase, is
\begin{equation}
 S_{mo}=\sqrt{\eta_o\eta_m}\,
       \frac{\texttt{ConverterEfficQ}}{\texttt{Ein}},\qquad
 \eta=\eta_o\eta_m\,\texttt{NormalizedEfficiencyQ},
 \label{eq:notebook-normalization}
\end{equation}
Here $\epsilon=1$ and
$\texttt{NormalizedEfficiencyQ}=\texttt{EfficiencyQ}/\texttt{Ein}^{2}$ for real
\texttt{Ein}. Examples set \texttt{Ein}=1; normalized linear efficiency is
independent of drive amplitude.

\texttt{SigmaPerp} implements Eq.~\eqref{eq:sigma};
\texttt{OnResonanceEfficiencyQ} implements Eq.~\eqref{eq:eta0} with unit port
fractions and $\gadd=\gamma_r$.

\texttt{parmsScreenshot} gives Eq.~\eqref{eq:benchmark}, while
\texttt{parmsOriginal} retains the original notebook's different parameters.
Keep $Q$ symbolic until applying \texttt{qRule}. The notebook's $x,z$ labels denote
the same two matched axes called $x,y$ here. Its \texttt{etaOpt},\texttt{etaMM}
are $\eta_o,\eta_m$, applied by \texttt{EfficiencyQWithPorts} at fixed linewidths.
That function multiplies the unnormalized \texttt{EfficiencyQ}; use
\texttt{Ein=1} or divide its result by $\texttt{Ein}^2$ for real nonzero drive.

\begin{thebibliography}{4}
\bibitem{bloch} I. Bloch, J. Dalibard and W. Zwerger, \emph{Rev. Mod. Phys.} \textbf{80}, 885 (2008),
 Sec. II. \href{https://doi.org/10.1103/RevModPhys.80.885}{doi:10.1103/RevModPhys.80.885}.
\bibitem{kumar} A. Kumar et al., \emph{Nature} \textbf{615}, 614--619 (2023).
 \href{https://doi.org/10.1038/s41586-023-05740-2}{doi:10.1038/s41586-023-05740-2}.
\bibitem{kumarSI} A. Kumar et al., \emph{Quantum-limited millimeter wave to optical
 transduction}, arXiv:2207.10121 (2022), Supplement B.1--B.3.
 \href{https://arxiv.org/html/2207.10121}{Full text and supplement}.
\bibitem{model} {\raggedright Baseline: \texttt{Abhishek\_Converter.nb}; QPM implementation:
 \texttt{Abhishek\_Converter\_QPM.nb}. Code and calculator:
 \href{https://github.com/AVKGitHub1/transduction_calcs}{AVKGitHub1/transduction\_calcs},
 commit \texttt{30da8841b3c9}. The $Q$ extension is derived here.\par}
\end{thebibliography}
\end{document}
